\clearpage
\section*{September 10, 2026: checking the construction-to-estimation handoff}
\addcontentsline{toc}{section}{September 10, 2026: checking the construction-to-estimation handoff}
\textit{Agent-drafted entry; main estimation inputs connected using preserved zoning.}

Checking the completed construction data against earlier recorded decisions found instructions affecting 109 projects that had not carried forward. The completed 268-case review covered a fixed subset, so it did not establish that every older decision was applied. The missing instructions have been restored to the existing cleaning stages. Newer explicit decisions take priority. The regression preparation code does not correct years, home counts, floor areas or land areas.

Relative to the 13,764-project dataset preserved at the start of this handoff, the rebuild removes 50 previously excluded records and changes construction years for 51 retained projects and home counts for five. No retained project's reported floor or land area changes. Six recorded decisions withhold density while preserving the buildings and their evidence. These counts describe the full construction dataset, not changes in the paper's 500-foot estimation sample. Restored year corrections also require locations and boundary distances to use the corrected year. The three location follow-ups now resolve through the existing source rules, and an unchanged second construction build requires no work.

Regression preparation attaches the existing boundary segments, daily alderman assignments, scores and ward controls to finished buildings, preserving the June 15 date proxy. Floor-area-ratio regressions now require usable floor and land area; homes-per-acre regressions require usable home counts and land area. A missing home count alone no longer excludes an otherwise usable floor-area ratio. An isolated main-regression run using the newly connected data succeeds. It uses 4,060 floor-area-ratio observations and 4,071 homes-per-acre observations, compared with the former common sample of 3,692. The corresponding multifamily counts rise from 822 to 863 and 874. At the handoff, production estimates had not yet been rerun.

Jacob approved using the preserved historical zoning table and reviewed map for this run. Zoning is assigned to all 9,530 records within 1,500 feet in the connected analysis data. Nine exact project/year matches use existing corrected-year zoning decisions. Reconstructing the preserved history from the original ordinances remains outstanding for raw-source replication.

Within 500 feet, the floor-area-ratio comparison retains 3,631 old observations, removes 61 and adds 429. The homes-per-acre comparison retains 3,632, removes 60 and adds 439. These are changes in sample records; changed identifiers and replacements need not represent newly discovered buildings. Of the retained observations, 99 in the floor-area-ratio sample and 100 in the homes-per-acre sample change at least one construction year, home count, floor area or land area. The old paper input remains preserved for comparison. The connected analysis input now replaces the old input, but paper estimates remained unchanged at that stage.

The main input build and targeted comparison reports pass. Specialized boundary-characteristic and project-to-permit inputs still require reconnection before a full paper rerun. The construction graph regenerates; the full paper graph still reaches an absent legacy task.

\textit{Reproduction note.} Branch: \texttt{sales\_sample\_exploration}. Run Make in the construction cleaning and analysis-data tasks. The release audit owns the two \texttt{construction\_estimation} comparison CSVs and preserves the old analysis input under \texttt{reference/}. The main-estimator check ran in a temporary copy.

\subsection*{Main density estimates rerun}
Jacob then requested the main density rerun. The production Make build succeeds and replaces the main density figure and sample summary. The estimated log difference on the more stringent side changes from $-0.131$ to $-0.104$ (standard error $0.047$) for floor-area ratio and from $-0.174$ to $-0.152$ ($0.073$) for homes per acre. For multifamily construction, the respective estimates change from $-0.149$ to $-0.138$ ($0.070$) and from $-0.194$ to $-0.198$ ($0.109$). All four remain negative. The overall estimates are smaller in magnitude; the multifamily estimates are less precise. The sample counts match the comparison above. These are the main specifications only, not a completed rerun of the robustness checks.

Jacob requested that the manuscript remain unchanged while he reviews these results. The temporary manuscript edits were reversed and the paper build was stopped before LaTeX compilation; the new main regression outputs remain available separately.

\subsection*{What accounts for the smaller overall estimates?}
A sequential decomposition reproduces both the old and new main estimates. On the project IDs common to the two samples, after updating measurements, years, locations and controls, the coefficients are $-0.132$ for FAR and $-0.180$ for DUPAC. Adding the 368 entering observations classified as individual houses or townhomes moves them to $-0.104$ and $-0.149$. Adding the entering multifamily observations then gives the final $-0.104$ and $-0.152$. Thus the added house and townhome observations account for almost all of the attenuation in the overall estimates in this ordering. Removing departing records makes both estimates more negative; corrections to retained records have a smaller net contribution. This is an accounting exercise whose individual contributions can depend on ordering, not a causal explanation of differences between building types.
The multifamily pattern differs: FAR attenuation comes mainly from measurement corrections among retained projects, while multifamily DUPAC does not attenuate overall. The audit's \texttt{density\_change\_decomposition.csv} records every fitted stage, sample size, estimate and standard error. The manuscript is unchanged.

\subsection*{Construction-year zoning review}
The 18 added sample records requiring a closer zoning check contain one confirmed broad-group error. The 2012 building at 1612 W Ontario inherited single-family zoning from a nearby preserved 2011 record. Ordinance O2011-8884 passed on December 14, 2011, changing the site to RT4 before the construction-date proxy. Its construction-year classification is multifamily. The other 17 retain their groups: ten have agreement across all five recorded maps, six have passed zoning amendments before construction consistent with the current map, and the Throop property has the same parcel in the preserved reviewed planned-development history and agreeing later maps.

The general correction prevents borrowing zoning from an adjacent year when the current map records an amendment between that year and construction. Across all 9,530 records, this changes four broad groups, including Ontario; three are within 500 feet. The other changes concern two Montana properties and one Washtenaw property. No buildings, construction dates, density measurements or locations change. The connected analysis data have been rebuilt; the regression estimates reported above have not been refreshed for these zoning-control corrections. This review uses the approved preserved history and does not establish full reconstruction of every historical ordinance.

\textit{Reproduction note.} The release audit's \texttt{added\_zoning\_eighteen\_review.csv} records the original cohort, map classifications, preserved historical support and passed ordinance evidence from the recorded City matters extract. Build its report target through the release-audit Makefile. The date check belongs in the construction zoning producer, rather than an observation-specific correction table. The manuscript remains unchanged.

\subsection*{Connecting the remaining density checks}
The zoning corrections leave the main FAR estimate at $-0.104$ and move DUPAC from $-0.152$ to $-0.151$, rounded to three decimals. The multifamily estimates remain $-0.138$ and $-0.198$. Sample counts remain 4,060 and 4,071 overall, and 863 and 874 for multifamily construction.

The boundary-characteristic task now recalculates its inputs for the connected construction dataset. Its road, water, park and cemetery measurements reproduce the old measurements across all 2,435 boundary segments up to numerical rounding. Among 3,645 retained project IDs in the old characteristics file, six change the straight-boundary indicator; the overlap indicators change for three of those six. The latter three have changed assigned segments, including one with a changed construction year. Amenity distances for most retained projects agree up to rounding. The location and boundary-restriction regression tables rebuild successfully.

Jacob approved retaining the recorded project-permit matches and adding exact-parcel or inside-parcel matches under the existing application and issue windows. Nearby permits alone do not qualify. A new producer records each match's source before the leave-project-out score check. This procedure does not change construction years or density measurements.

The clean-checkout test exposed an absent parcel-address source task and required historical source files not distributed in Git. The address task now restores a checksum-verified recorded snapshot. Jacob approved uploading the 103-file recorded-source archive to the repository's GitHub release. Its compressed size is 968,577,788 bytes. The source bootstrap has passed download-checksum verification, missing-file restoration and unchanged-build checks on GNU Make 3.81. Full construction reconstruction in the isolated checkout is still being checked; these packaging checks alone do not establish full paper replication.

\subsection*{Added observations are concentrated in the historical second ward}
The saved attenuation comparison adds the same 368 observations outside the
external multifamily category to both the FAR and DUPAC samples. Of these, 355
report one dwelling and 13 report multiple dwellings. Thus the earlier description
of all 368 as individual homes or townhomes was too broad.

The historical second ward accounts for 113 additions: 77 under Robert Fioretti,
35 under Madeline Haithcock, and one under Brian Hopkins. Of these, 106 have
construction years 2006 or 2007. This predominantly concerns the old ward geography.
Other concentrations are ward 26 (31), ward 35 (26), and wards 1 and 32 (22 each).
There are 218 additions on the lower-score side and 150 on the higher-score side;
111 of the second ward's 113 additions are on the lower-score side.
All recorded own-ward and neighboring-ward aldermen match the recorded terms at
the construction-date proxy used in the analysis. This check does not independently
reverify the locations or establish any ward's contribution to attenuation.

The audit reads the saved comparison and existing analysis, terms and score
files, without rebuilding production. Its output is
\texttt{entering\_home\_assignments.csv} in the working-paper release audit.
The task README records the Make invocation that holds those inputs fixed.
This is the first report for this 368-row assignment audit.

Jacob paused the rebuild and requires all necessary files to be committed before
the next clean-clone test. A missing committed input or a build failure requires
a repository fix and a new clone, rather than patching and continuing the failed
clone. The stopped exploratory run is not the final replication test. The
manuscript remains unchanged.

\subsection*{Historical assessments contradict some selected construction years}
A follow-up review of all 113 second-ward additions found a substantive problem.
All selected years match their cited Assessor rows, but 37 properties already
have the same reported floor and land areas in an assessment before the selected
construction year. Thirty-five are selected as 2006 and two as 2007. For example,
a home at 26 South Aberdeen has 1,778 square feet of building and 1,106 square feet
of land in 2005. Its reported construction year is 2004 through the 2021
assessment, then changes to 2006 in 2022 without a change in those measurements.
Eleven additional properties selected as 2007 have earlier reports of 2005;
their available histories begin in 2009, so the contradiction is less direct.

These dates were not manually invented. Selecting the later Assessor report
allowed conflicting year revisions into the analysis. The 37 direct conflicts
cannot be treated as verified construction dates. A general chronological check
is needed before the production rerun, with care not to confuse an older building
with a replacement building on the same parcel. No production years or sample
inclusions have been changed by this audit. The remaining properties' identities
and shared-land coverage are not certified by this year check.

The saved audit is \texttt{ward\_two\_year\_review.csv}; the release-audit README
records its frozen-input Make invocation. Its first standard report contains
113 unique project records. The rebuild remains paused.

\subsection*{Earlier Assessor years are not uniformly better}
Web research and preserved permits distinguish the 37 direct chronological
conflicts from the 11 Union Row homes. A purchase recorded in March 2004 at
26 South Aberdeen corroborates early development there, although the page does
not identify the individual unit. Current listing years disagree and do not
resolve the chronology by themselves. Selecting the latest reported year no
later than the first assessment with the same measurements would give 33 dates
of 2004 and four of 2005 for the 37 conflicts, putting them outside the study.
This is a proposed correction, not an applied result.

For Union Row, original new-building permits issued May 17, 2007 explicitly
cover the affected addresses. Thus automatically replacing 2007 with the older
Assessor report of 2005 would be inappropriate. Belgravia dates completion of the
whole 35-home development to 2010, without identifying individual completion
dates. The individual Assessor year 2007 remains a proxy, not a independently
verified completion date. The proposed general rule is to select a reported
Assessor year consistent with the identified building's assessment and original
permit history, keeping contradictory cases unresolved rather than inventing a
year. The audit README records URLs and limitations. Production and estimation
remain unchanged.

\subsection*{Population trial of chronological year selection}
Jacob approved correcting the 37 and requested a population check before broader
adoption. The existing assessment-selection script now tests the chronological
rule. All 37 receive the approved pre-period dates: 33 in 2004 and four in 2005.
The eleven Union Row dates remain 2007. Across 28,100 residential candidates,
981 years change; units, floor area, land area and measurement-row identifiers
are unchanged. The build used preserved inputs and stopped at selection.

Against the saved final population, 869 projects would change, including 273
within 500 feet. Sixty of those 273 would move before 2006. Twenty-six proposed
changes exceed three years and require scrutiny before broad adoption. These
are trial results, not adopted changes to the final sample or new estimates.

A separate screen covers all 13,714 final projects using the preserved residential,
commercial and condominium measurements. It finds 924 potential contradictions,
11,958 without a direct contradiction and 832 without a comparable historical
snapshot. The latter are unverified by this check, not identified errors.
The commercial contradictions concern 4654 North Sheridan and 202 West Hill:
2021 source records report 2020 construction, while later records report 2022.
No commercial year change has been applied. Shared sites and changing card sets
need more than measurement equality to establish building continuity.

The population audit and its standard report record every project's status;
the release-audit README gives the frozen-input Make command. Final dataset,
geography, estimation and manuscript remain unchanged. The draft rule requires
review before an ordinary downstream build is resumed.

\subsection*{Review of all 26 large year changes}
The bounded review recommends retaining seven as new construction, excluding
fourteen existing buildings or additions, excluding two pre-2006 buildings, and
holding out three whose new-building identity is insufficiently supported.
These remain recommendations. Only four are in the current 500-foot sample:
retain 2425 West Arthington, exclude 2340 North Leavitt and 5921 South Monitor,
and hold out 1308 North La Salle.

The draft rule would incorrectly assign some replacements their demolished
building's year. Arthington has original 2008 demolition and new-building permits;
a listing explicitly describes a 2010 rebuild. Early assessments nevertheless
carry 1892 onto the new building's measurements. Exact-address searches also found
new-building permits missing from existing project links. The chronological rule
must account for those original permits and building episodes before selecting
an older year. The review does not justify adopting all draft corrections.

The saved 26-row review provides every address, proposed treatment, evidence and
remaining uncertainty. Overhill's 2008 and Ridgeway's 2008 remain Assessor proxies,
not independently established completion dates. The production dataset and
manuscript remain unchanged; the release-audit README records reproduction.

\subsection*{Approved correction limited to 37 buildings}
The broad automatic construction-year rule is withdrawn. Jacob approved correcting
only the 37 reviewed buildings whose matching historical Assessor records support
pre-period construction. Thirty-three receive 2004 and four receive 2005. The
existing recorded-year input applies these decisions before project selection;
these buildings therefore fall outside the 2006--2022 analysis period. The other
recommendations from the 26-case review remain unadopted.

Comparison with the assessment selection saved before the experiment confirms
exactly these 37 construction-year changes. Dwelling counts, floor areas and land
areas are unchanged throughout that comparison. This incremental rebuild does
not replace the planned clean-clone replication test. The manuscript remains
unchanged.


The rebuilt construction ledger has 13,677 records, down from 13,714; the analysis
data have 9,493, down from 9,530. Every surviving record is unchanged in every
column. All 37 removed observations were within 500 feet and outside the
multifamily subsample. The all-construction FAR estimate moves from $-0.104$
(SE 0.047) to $-0.097$ (SE 0.045); DUPAC moves from $-0.151$ (SE 0.073) to
$-0.148$ (SE 0.073). Both remain significant at five percent. Their regression
samples contain 4,023 and 4,034 observations. Multifamily estimates and sample
counts are unchanged. These corrections do not reverse the previously observed
attenuation.

The incremental build used Make in the density main-results task, followed by
the producing tasks' data reports. A subsequent dry run schedules no substantive
producer. The saved density figure was inspected; the manuscript was not compiled. An accidental root Make invocation was stopped
during upstream checks before document compilation.

\subsection*{Score sensitivity rechecked before the release review}
On the preserved permit input, sequential and joint estimation each use 120,505
permits. The correlation between their unshrunk alderman effects remains 0.348
(rank correlation 0.339). The joint model changes the control coefficients as
well as the alderman effects: the income coefficient per \$10,000 changes from
$-0.0150$ to $-0.0626$. This comparison does not establish that either score
identifies a causal alderman effect. Neighborhood characteristics and alderman
identity overlap substantially by construction; the two procedures allocate
these differences differently.

The absorbed joint fit warns about redundant fixed effects. Re-estimating with
explicit alderman indicators, using a common reference alderman, reproduces its
alderman differences within $5.3\times10^{-7}$ log points. This rules out that
normalization issue as the explanation for the observed ranking change. The
comparison is recorded in the existing score audit, including the control
coefficients and an assertion checking equivalent parameterizations. Production
scores and the manuscript have not been changed.

The broader production review is ongoing. Known open questions include the
score interpretation, the incomplete first-stage exclusion in the own-project
score sensitivity, rental date assignment and aggregation, and live-source
vintages outside the preserved construction archive. A checkpoint commit should
not be read as clearance of these issues or completion of the fresh-clone test.

\subsection*{Decision to retain the current score}
Jacob chose to retain the existing sequential score for the conference draft.
Its intended interpretation is alderman behavior after adjustment for common
neighborhood and permit differences. The joint-estimation comparison remains
sensitivity evidence, not a replacement production specification. Remaining
delays may also contain unobserved neighborhood or project differences; the
score is not a separately identified causal alderman effect. This decision
resolves which score the forthcoming rerun should use, without changing the
construction corrections or clearing the other open review items.

\subsection*{Fresh-clone download failure}
The conference replication attempt at commit \texttt{29dc42a4} reproduced all
13,677 construction records exactly and the main density coefficients at their
reported precision. It did not complete: one of 470 Cook County sales-location
requests failed after three attempts. The identical request subsequently returned
all 500 parcel records in approximately ten seconds, supporting a temporary
service failure rather than unavailable data. This run is not a successful
end-to-end replication.

The downloader now retries unsuccessful requests after the other batches finish,
giving the service time to recover. Successful requests are not repeated, and
an unresolved failure still stops the build before publishing data. A disposable
test verified both recovery from a temporary failure and rejection of a persistent
failure. The full build must restart from a new remote clone with an empty Census
geography cache. No failed-clone outputs will be supplied to that test.

Jacob subsequently revised the testing sequence: complete the ongoing build as
a diagnostic run, fixing and resuming after failures instead of restarting each
time. Collect and commit the repairs before one final fresh-clone validation.
A successfully patched diagnostic build will not be described as an unattended
replication.

The diagnostic run exposed a further download risk: batches contained 24 requests,
whereas the installed curl pool permits six connections per host. Curl counts
queueing time toward connection timeouts, so waiting requests can expire before
they connect. The batch size was reduced to six without changing query contents.
All 24 requests in a live four-batch check succeeded. This removes the unnecessary
queue; it does not guarantee that the external service will always respond.

The six-request full download still left four 2022 requests unsuccessful after
three attempts. A direct check of batches 9--12 returned successful responses
in 24--44 seconds. Connection and total time limits were increased to 120 and
300 seconds, respectively, with at most five attempts and explicit error
messages identifying unsuccessful requests. The diagnostic build was resumed
again, retaining completed upstream analyses.

Two downstream implementation issues were also corrected during the diagnostic
run. The price-results consistency check now computes the range of original
distances before replacing them with their median; previously it always found
a zero range. The same task now regenerates its jointly produced figures when
one is missing, declares its Makefile and package setup as dependencies, and
produces a standard report for its saved coefficient table. These changes do
not alter the regression specification.

\subsection*{Diagnostic build completed}
The diagnostic checkout compiled the main paper (31 pages), online appendix
(21 pages), combined working paper (56 pages), and word-count document
(20 pages), with no unresolved references or LaTeX errors. Its construction
ledger exactly matches the approved 13,677-record file. The 231,800 sales-parcel
coordinate records also match the previous file byte for byte. With six-request
batches and the longer time limits, all 470 source requests succeeded on the
first attempt.

The main full-sample density coefficients remain $-0.097$ for FAR and $-0.148$
for DUPAC at reported precision. The main rent and sales coefficients are
approximately $0.0175$ and $0.0300$, matching the prior estimates to numerical
precision, with unchanged samples of 240,713 rental records and 58,468 sales.
The refreshed live permit source produces small score differences without
changing these main results or boundary-side assignments. The repaired distance
check found no contradiction in the prior inputs: the maximum within-location
difference was less than $4\times10^{-9}$ feet across 81,315 locations.

All working-paper pages were rendered and inspected. One density table extends
slightly beyond the text width without clipping. Manuscript prose remains
unchanged at Jacob's request and still contains older density estimates and
sample counts; compilation alone does not make it ready to circulate. The
patched diagnostic build is not the final fresh-clone validation.

GNU Make 3.81 regenerated missing price figures with one producer invocation
under a four-job build and refreshed the coefficient report in that same run.
Deleting only the parcel report regenerated only that report, without another
download. An unchanged price-task build and an unchanged full paper build ran
no substantive producer or document compiler. These checks completed in the
diagnostic checkout before the fixes were prepared for the final fresh clone.
